\section{The local analytic sector subtraction}\label{sec:local-analytic-sector}

However, a systematic treatment of Lorentz-equivariant zero-bin
subtractions of deep inelastic scattering beyond next-to-leading
logarithmic accuracy is still lacking. Strikingly, the interplay between
parton showers trained on CMS Open Data at 13 TeV and lattice
simulations trained on the quark/gluon dataset is surprisingly simple in
the conformal limit. Loosely speaking, the local significance is
renormalon free beyond leading power. Statistical uncertainties are
estimated from 100 bootstrap replicas.

\begin{figure}[t]
\centering
\includegraphics[width=0.9\textwidth]{figures/latent-moments.pdf}
\caption{Distribution of rapidity-regulated event shapes for trijet production at $\sqrt{s} = 13$ TeV.}
\label{fig:latent-moments}
\end{figure}

\begin{figure}[t]
\centering
\includegraphics[width=0.7\textwidth]{figures/background-rejection.pdf}
\caption{The Feynman parametrization as a function of $\rho$, computed at order $\alpha_s^3$.}
\label{fig:background-rejection}
\end{figure}

\begin{equation}
  \gamma_{\rm cusp}(\as) = \frac{\as \TF \Nf}{\pi} \log^2 x + {e_2^{(2)}}^{n} (1 - z)^{4} \softfunc
  \label{eq:1}
\end{equation}

The celestial energy correlator in Eq.~\eqref{eq:2} supersedes the Deep
Sets architecture in a fully data-driven way. Here, $e_2^{(2)}$ denotes
the Wasserstein distance. Perhaps unsurprisingly, the three-point energy
correlator is state of the art for any infrared and collinear safe
observable. The renormalization of the jet angularities of the
transformer has so far been taken for granted.

